\documentclass[10pt,a4paper]{amsart}
\usepackage[margin=19mm]{geometry}
\usepackage{amsmath,amssymb,mathtools}
\usepackage[scr=rsfs,cal=euler]{mathalfa}
\usepackage{xfrac}
\usepackage[unicode,hidelinks,bookmarks=false]{hyperref}
\newcommand{\ie}{i.e.\ }
\newcommand{\cf}{cf.\ }
\newcommand{\resp}{respectively\ }
\newcommand{\Subsection}{\subsection}
\providecommand{\nobreakdash}{\nobreak\hbox{-}}
\setlength{\emergencystretch}{2em}
\multlinegap0pt
\usepackage{bm}
\usepackage{bbm}
\usepackage{dsfont}
\usepackage{xspace}
\usepackage{cancel}
\usepackage{mathtools}
\usepackage{amsthm}
\makeatletter
\@ifundefined{Thm}{\newtheorem{Thm}{Thm}}{}
\@ifundefined{Lemma}{\newtheorem{Lemma}[Thm]{\sc Lemma}}{}
\@ifundefined{Theorem}{\newtheorem{Theorem}[Thm]{\sc Theorem}}{}
\makeatother
\newcommand{\QQ}{{\mathbb Q}}
\newcommand{\Refl}{{\mathop{\operatorname{Ref}\, }}}
\newcommand{\ZZ}{{\mathbb Z}}
\newcommand{\cA}{{\mathcal A}}
\newcommand{\cE}{{\mathcal E}}
\newcommand{\cO}{{\mathcal O}}
\newcommand{\ch}{{\mathop{\rm ch \, }}}
\newcommand{\rk}{{\mathop{\rm rk \,}}}
\title[Intersection theory and Chern classes on normal varieties]{Intersection theory and Chern classes on normal varieties}
\author{Adrian Langer}
\date{Source: arXiv:2210.08766v4 (CC BY 4.0)}
\begin{document}
\maketitle

\noindent\emph{Source and licence.} Intersection theory and Chern classes on normal varieties. Version arXiv:2210.08766v4 (CC BY 4.0), distributed under the Creative Commons Attribution 4.0 International licence, as recorded in the arXiv licence metadata for this exact version and shown on the abstract page https://arxiv.org/abs/2210.08766v4. Full rights notice: licenses/arxiv-2210.08766-CC-BY-4.0.txt.\par\smallskip

\noindent\textit{Editorial scope.} Counted unit: the Theorem whose source label is \texttt{Bogomolov's-inequality} in the article (source0.tex lines 1677--1683, source label \texttt{Bogomolov's-inequality}), delivered together with the article's own complete original proof of it (source0.tex lines 1685--1709, one \texttt{proof} environment of 25 lines). No mathematics is written, repaired or re-derived: statement and proof are the article's own text line for line. Required context retained uncounted: finite-generation (lines 649--651); main-consequence (lines 1395--1404); boundedness-3 (lines 1654--1660). Every definition, display and label of those lines is retained unchanged. Omitted from this package and not redistributed: 1--648, 652--1394, 1405--1653, 1661--1676, 1684--1684, 1710--1942. Nothing inside a retained excerpt is omitted or altered: every retained body is delivered exactly as the article's lines are written, so no omission receipt of any kind accompanies this package. Citations: the retained excerpts carry no citation command at all, so no bibliography entry of the article is transcribed and no reference is redistributed. Reference adaptation: none was needed and none was made; every reference inside the delivered unit resolves inside it, so no unresolved reference or citation is produced. Printed slips kept literal: no printed word, symbol, hypothesis or number of the counted statement or of its proof was altered. Layout and class: the article's own class is \texttt{article} with packages including a4wide, amsmath, amsthm, amssymb, amscd, mathptmx, amsfonts, color, pstricks, xy; the wrapper is the lane's pinned \texttt{amsart} editorial wrapper at 10pt on A4, which restates the theorem-like environments the retained text needs and adds the article's own macro definitions that the retained text needs (\texttt{\textbackslash QQ}, \texttt{\textbackslash Refl}, \texttt{\textbackslash ZZ}, \texttt{\textbackslash cA}, \texttt{\textbackslash cE}, \texttt{\textbackslash cO}, \texttt{\textbackslash ch}, \texttt{\textbackslash rk}), with extra wrapper packages (bm, bbm, dsfont, xspace, cancel, mathtools). The delivered numbering is the unit's own. Assets: no retained span contains a figure environment, a graphics-inclusion command, a PDF-page inclusion, a table or a TikZ picture; no image file is copied into this package, the package declares no asset dependency and no image file is redistributed with it. Licence: version arXiv:2210.08766v4 of this article is distributed under the Creative Commons Attribution 4.0 International licence, as recorded in the arXiv licence metadata for this exact version and shown on the abstract page https://arxiv.org/abs/2210.08766v4. A full rights notice is delivered at licenses/arxiv-2210.08766-CC-BY-4.0.txt. Characters: every retained excerpt is pure ASCII, so the delivered engine, XeLaTeX with its default Latin Modern text font, reports no missing character.
\begin{Lemma}\label{finite-generation}
$N_L(X)$ is a free $\ZZ$-module of finite rank. In particular, there exists a positive integer $N$ such that the intersection pairing $\langle \cdot , \cdot \rangle$ takes values in $\frac{1}{N}\ZZ\subset \QQ$.  If  $\rk N_L (X)=s$ then the intersection pairing $\langle \cdot , \cdot \rangle$ has signature $(1, s-1)$.
\end{Lemma} 

\begin{Theorem}\label{main-consequence}
Let  $L_1, ..., L_{n-2}$ be very ample line bundles on $X$. Then there exists a constant $C$ such that for all
$\cE\in \Refl(\cO_X)$ we have
\begin{align*}
&\left\lvert \chi (X, c_1(L_1)... c_1(L_{n-2}) \cdot \cE) -r\chi (X, c_1(L_1)... c_1(L_{n-2}) \cdot \cO_X) -  \int _X \ch _2 (\cE)L_1...L_{n-2}  \right. \\
 & \left. +\frac{1}{2}   c_1(\cE).(K_X+L_1+...+L_{n-2}). L_1...L_{n-2} 
\right\rvert \le C\cdot r^2,
    \end{align*}
where $r$ is the rank of $\cE$.
\end{Theorem}

\begin{Theorem} \label{boundedness-3}
Let us fix some positive integer $r$  and some real
numbers $c_2$ and $\mu _{\max}$. Let us also fix some class $c_1\in N_H(X)$.
Then the set $\cA$ of coherent reflexive $\cO_X$-modules $\cE$ of rank $r$ with $[c_1(\cE)]=c_1\in N_H(X)$,
$\int _X c_2 (\cE) H^{n-2}\le c_2$  and $\mu _{\max ,H } (\cE)\le \mu _{\max}$
is bounded.
\end{Theorem}

\begin{Theorem}\label{Bogomolov's-inequality}
Let us fix some positive integer $r$ and some non-negative rational number $\alpha$.
There exists some constant $\tilde C=\tilde C(X,H,r,\alpha)$ depending only on $X$, $H$, $r$ and $\alpha$ such that for every coherent
reflexive    $\cO_X$-module $\cE$ of rank $r$ with
$\mu _{\max ,H } (\cE)-\mu _H (\cE)\le \alpha$ we have
$$\int _X \Delta (\cE) H^{n-2}\ge \tilde C.$$
\end{Theorem}

\begin{proof}
Let us choose a basis $L_1,....,L_s$ of $N_H (X)$ as a $\ZZ$-module (see Lemma \ref{finite-generation})
and let us write $[c_1 (\cE)]= \sum a_i L_i$ for some $a_i\in \ZZ$. There exist  uniquely determined integers $q_i$
and $r_i$ such that $a_i=q_i r+r_i$ and $0\le r_i <r$.
Since 
$$\int _X \Delta (\cE) H^{n-2}= \int _X \Delta (\cE (-\sum q_i L_i)) H^{n-2}$$
and there are only finitely many possibilities for  $[c_1 (\cE (-\sum q_i L_i))]= \sum r_i L_i$, it is sufficient to prove existence of the above constant assuming that $c_1=[c_1(\cE)]\in N_H (X)$ is fixed.

If $\int _X c_2 (\cE) H^{n-2}\ge 0$ then  $\int _X \Delta (\cE) H^{n-2}\ge - (r-1)c_1^2.H^{n-2}$.

On the other hand, by Theorem \ref{boundedness-3}  the set $\cA$ of reflexive coherent $\cO_X$-modules $\cE$ of rank $r$ with $[c_1(\cE)]=c_1\in N_H(X)$, $\int _X c_2 (\cE) H^{n-2}\le 0$  and $$\mu _{\max ,H } (\cE)\le \mu _{\max}= \alpha+\frac{1}{r}c_1.H^{n-1}$$ is bounded. So there exists some constant $D$ such that for every $\cE\in \cA$ we have
and any $\cE$-regular sequence $H_1,..., H_{n-2}\in |\cO_X (1)|$ we have
$$\chi (X, H^{n-2}\cdot \cE)=\chi (X, \cE |_{\bigcap _{j\le n-2} H_j })\le D.$$
But Theorem \ref{main-consequence} gives
\begin{align*}
\chi (X, H^{n-2} \cdot \cE) \ge &\frac{1}{2r}   c_1.(c_1-r(K_X+(n-2)H)). H^{n-2} 
 -\frac{1}{2r}\int _X \Delta (\cE)H^{n-2}   \\
 &  +r\chi (X, H^{n-2} \cdot \cO_X) - C\cdot r^2.
    \end{align*}
Therefore for any $\cE\in \cA$ we have
\begin{align*}
\int _X \Delta (\cE) H^{n-2}\ge \max &(- (r-1)c_1^2.H^{n-2} ,c_1 . (c_1 -r(K_X+(n-2)H).H^{n-2}\\
&+ 2r^2\chi (X, H^{n-2} \cdot \cO_X )-2C\cdot r^3 +2D\cdot r ).
\end{align*}
\end{proof}

\end{document}
